\documentclass[aspectratio=169]{beamer}
\usepackage{will_handley_beamer}
\usepackage{title_page}

% Commands
% --------
% - \arxiv{arxiv number}
% - \arxiv{<number>}            arxiv.org/abs/<number>
% - \oldarxiv{<arxiv number>}   arxiv.org/<number>
% - \doi{<doi>}                 doi.org/<doi>
% - \xkcd{<number>}             xkcd.com/<number>
% - \email{<email>}             <<email>>
% - \tthref{<website>}          <website>
% - \av[dist]{<quantity>}       <quantity>_{dist}
% - \student{<name>}{<detail>}{<photo>}

% Talk details
% ------------
\title{AI/ML Tools for Research}
\subtitle{KICC Workshop}
\date{23rd July 2025}

\begin{document}

\begin{frame}
    \titlepage
\end{frame}

\begin{frame}
    \frametitle{Overview}
    
    \begin{columns}
        \column{0.5\textwidth}
        \begin{itemize}
            \item Follow-up to successful Part II/III training.
            \item Focus on ``doing'' not ``explaining''.
            \item Three layers of AI tools framework.
            \item Interactive demonstrations.
            \item 90 minutes then tea.
        \end{itemize}
        
        \vspace{0.3cm}
        
        \textbf{Common questions:}
        \begin{itemize}
            \item ``Which model should I use?''
            \item ``Don't these things hallucinate?''
            \item ``How do I get started?''
            \item ``Is this worth the investment?''
        \end{itemize}
        
        \column{0.5\textwidth}
        \begin{block}{Previous session covered:}
            \begin{itemize}
                \item Academic practice and Cambridge policies.
                \item How LLMs work and their limitations.
                \item Bias, hallucinations, and mitigation.
                \item Privacy and intellectual property.
                \item Model landscape and responsible usage.
            \end{itemize}
        \end{block}
    \end{columns}
    
    \vfill
    \centering
    \footnotesize{Previous slides: \tthref{docs.google.com/presentation/d/1rnbk7WZUjMIAGBzkN6i8v0tEe1imiHprvL1ZOjs2B20}}
\end{frame}

\begin{frame}
    \frametitle{First things first}
    If you've been busy, the past six months have seen a couple of inflection points
    \vspace{0.5cm}
    \begin{itemize}
        \item \textbf{February 2025}: Models acquired many of the capabilities we prize in our PhD students
            \begin{itemize}
                \item Models: o3 series, Gemini 2.5, Claude 3.5 Sonnet.
                \item Capabilities: Code development, mathematical reasoning, literature review, paper \& grant drafting.
            \end{itemize}
        \item \textbf{May 2025}: Agentic systems launched commercially
            \begin{itemize}
                \item Tools: Claude Code, Cursor agent mode, Deep Research.
                \item Capabilities: Writing test suites, assembling pip-installable codes, synthesising handwritten notes and TeX files, write/run/debug code in languages you don't know.
            \end{itemize}
        \item \textbf{The performance gap}: Free models lag significantly behind paid versions
    \end{itemize}
    
    \vspace{0.5cm}
    If your lived experience is typing a research question into free ChatGPT at the start of the year and finding it was rubbish, things have moved on
\end{frame}

\begin{frame}
    \frametitle{On AI Hype}
    \begin{columns}
        \column{0.48\textwidth}
        \begin{itemize}
            \item These things are not ``intelligent'' in the way Silicon Valley wants you to believe (AGI is not just around the corner).
            \item \emph{Caveat emptor}: It is very hard for a human to read language, and not reflexively construct a mind behind it.
            \item You should not however think of them as mere ``next token predictors'':
                \begin{itemize}
                    \item The modern reinforcement learnt reasoning systems are significantly more than that.
                    \item Agentic tools ground these in reality.
                \end{itemize}
        \end{itemize}
        
        \column{0.48\textwidth}
        \begin{block}{A senior examiner's experiment}
            \begin{itemize}
                \item I got o3 to take the Part II exams.
                \item Prompt on exam day: ``You are a first class Part II Cambridge astronomy student... here is the syllabus... here is the question, answer it''.
                \item Toby Lovick transcribed answers onto written exam scripts.
                \item Slipped them amongst real scripts and marked blind.
                \item \textbf{Result}: Best student we've had across IoA history, even without coursework.
            \end{itemize}
        \end{block}
    \end{columns}
\end{frame}

\begin{frame}
    \frametitle{The AI landscape}
    \centering
    \begin{tabular}{l|c|c|c|}
        \visible<3->{Under the hood} & 
        \visible<3->{\includegraphics[valign=m, width=0.2\textwidth]{logos/microsoft}} & 
        \visible<3->{\includegraphics[valign=m, width=0.2\textwidth]{logos/amazon}} & 
        \visible<3->{\includegraphics[valign=m, width=0.2\textwidth]{logos/google}} \\[0.5cm]
        
        \visible<2->{Companies} & 
        \visible<2->{\includegraphics[valign=m, width=0.2\textwidth]{logos/openai}} & 
        \visible<2->{\includegraphics[valign=m, width=0.2\textwidth]{logos/anthropic}} & 
        \visible<2->{\includegraphics[valign=m, width=0.2\textwidth]{logos/deepmind}} \\[0.5cm]
        
        \visible<1->{Models} & 
        \visible<1->{\includegraphics[valign=m, width=0.2\textwidth]{figures/chatgpt.png}} & 
        \visible<1->{\includegraphics[valign=m, width=0.2\textwidth]{figures/claude.pdf}} & 
        \visible<1->{\includegraphics[valign=m, width=0.2\textwidth]{figures/gemini.pdf}}
    \end{tabular}
    
    \vspace{1cm}
    \visible<4->{
        \textbf{Others exist:} Perplexity, Poe, Character.ai, You.com, Cohere, xAI, Llama...
    }
    
    \vspace{0.5cm}
    \visible<5->{
        \textbf{But these are the three main pillars (July 2025).}
    }
\end{frame}

\begin{frame}
    \frametitle{Three layers of AI tools}
    \framesubtitle{A framework for understanding the landscape}
    \begin{block}{Layer 3: Agentic Systems}
        \vspace{-0.3cm}
        \begin{columns}
            \column{0.48\textwidth}
            \begin{itemize}
                \item Claude Code, cursor agent mode.
                \item Custom workflows.
            \end{itemize}
            \column{0.48\textwidth}
            \begin{itemize}
                \item Autonomous task completion.
                \item Steeper learning curve but transformative.
            \end{itemize}
        \end{columns}
    \end{block}
    
    \vspace{0.5cm}
    
    \begin{block}{Layer 2: Chat-based AI}
        \vspace{-0.3cm}
        \begin{columns}
            \column{0.48\textwidth}
            \begin{itemize}
                \item ChatGPT, Claude, Gemini.
                \item Web interfaces.
            \end{itemize}
            \column{0.48\textwidth}
            \begin{itemize}
                \item Interactive problem solving.
                \item Good for exploration and learning.
            \end{itemize}
        \end{columns}
    \end{block}
    
    \vspace{0.5cm}
    
    \begin{block}{Layer 1: Autocomplete}
        \vspace{-0.3cm}
        \begin{columns}
            \column{0.48\textwidth}
            \begin{itemize}
                \item GitHub Copilot.
                \item VS Code extensions.
            \end{itemize}
            \column{0.48\textwidth}
            \begin{itemize}
                \item Completes code as you type.
                \item Minimal learning curve.
            \end{itemize}
        \end{columns}
    \end{block}
\end{frame}

\begin{frame}
    \frametitle{Layer 1: Autocomplete}
    \framesubtitle{Standard code-completion, but powered by AI}
    
    \begin{columns}
        \column{0.6\textwidth}
        \begin{block}{Core idea}
            \begin{itemize}
                \item AI-powered code-completion.
                \item Trained on all of GitHub.
                \item Context-aware suggestions.
            \end{itemize}
        \end{block}
        
        \begin{block}{Recommendation}
            \begin{itemize}
                \item GitHub Copilot (Pro subscription).
                \item Free for university email holders.
                \item \url{github.com/settings/education/benefits}.
                \item Worth \$10/month.
                \item Also available in Cursor.
            \end{itemize}
        \end{block}

        \column{0.4\textwidth}
        \includegraphics[width=\textwidth]{figures/copilot_screenshot.png}
    \end{columns}
\end{frame}

\begin{frame}
    \frametitle{Layer 2: Chat-based AI}
    \framesubtitle{The usual chat versions}
    
    \begin{columns}
        \column{0.3\textwidth}
        \centering
        \includegraphics[width=\textwidth]{figures/chatgpt_screenshot.png}
        
        ChatGPT
        
        \url{chat.openai.com}
        
        \column{0.3\textwidth}
        \centering
        \includegraphics[width=\textwidth]{figures/claude_screenshot.png}
        
        Claude
        
        \url{claude.ai}
        
        \column{0.3\textwidth}
        \centering
        \includegraphics[width=\textwidth]{figures/gemini_screenshot.png}
        
        Gemini
        
        \url{gemini.google.com}
    \end{columns}
\end{frame}

\begin{frame}
    \frametitle{Layer 2: Developer versions}
    \framesubtitle{Use their developer level tools}
    
    \begin{columns}
        \column{0.3\textwidth}
        \centering
        \includegraphics[width=\textwidth]{figures/chatgpt_dev_screenshot.png}
        
        OpenAI Playground
        
        \url{platform.openai.com/playground}
        
        \column{0.3\textwidth}
        \centering
        \includegraphics[width=\textwidth]{figures/claude_dev_screenshot.png}
        
        Claude Workbench
        
        \url{console.anthropic.com/workbench}
        
        \column{0.3\textwidth}
        \centering
        \includegraphics[width=\textwidth]{figures/gemini_dev_screenshot.png}
        
        Google AI Studio
        
        \url{aistudio.google.com}
    \end{columns}
    
    \vspace{0.3cm}
    
    \textbf{These should be your default as a scientist} -- if you're happy using a command line, you shouldn't be satisfied by an interface that looks like a child's toy.
\end{frame}

\begin{frame}
    \frametitle{Layer 2: Chat-based AI}
    \framesubtitle{Use AI to improve AI}
    \begin{columns}
        \column{0.48\textwidth}
        \begin{itemize}
            \item Likely the layer of AI most familiar to you.
            \item ``chatting''/conversation can be very powerful for naturally tuning the attention/context of the model.
            \item For one shot work, a little prompt engineering goes a long way:
                \begin{itemize}
                    \item You can use these tools to improve prompts.
                \end{itemize}
        \end{itemize}
        \column{0.48\textwidth}
        \includegraphics<1>[width=\textwidth]{figures/prompt_demo1.png}%
        \includegraphics<2>[width=\textwidth]{figures/prompt_demo2.png}
    \end{columns}
\end{frame}

\begin{frame}
    \frametitle{Layer 2: Cost}
    \framesubtitle{Pay as you use vs subscription models}
    
    \begin{columns}
        \column{0.6\textwidth}
        \begin{itemize}
            \item Gemini is free to use (but you sign away your data).
            \item ChatGPT and Claude offer \$20/month subscriptions.
            \item API access available: pay per token consumed:
                \begin{itemize}
                    \item \$0.0002-0.015 per 1000 tokens (roughly 750 words).
                    \item Orders of magnitude cheaper for research use.
                    \item Only pay for what you actually use.
                \end{itemize}
            \item Never buy yearly subscriptions - AI evolves too rapidly.
            \item My group's preference: pay for API access via tokens.
                \begin{itemize}
                    \item Pay for what you use.
                    \item Allows you to try out the latest models.
                    \item API keys allow you to do all kinds of tricks at the command line.
                \end{itemize}
        \end{itemize}
        
        \column{0.4\textwidth}
        \includegraphics[width=\textwidth]{figures/cost_demo1.png}
        \includegraphics[width=\textwidth]{figures/cost_demo2.png}
    \end{columns}
\end{frame}

\begin{frame}
    \frametitle{Layer 3: Agentic Systems}
    \framesubtitle{Autonomous task completion - transformative but complex}
    
    \begin{columns}
        \column{0.6\textwidth}
        \begin{itemize}
            \item Emerged commercially in April 2025.
            \item If you get too skilled at chat-based AI, your job ends up being a copy-paste-run-debug monkey.
            \item Agentic systems are what happens if you let ChatGPT edit files, run commands, and inspect the behaviour.
            \item This closes the scientific loop.
            \item Agentic systems are (partially) grounded in reality. Correctly guided, they can reduce hallucinations by testing their hypotheses.
            \item Claude Code (CLI), Cursor Agent mode, Gemini CLI.
            \item Deep research also falls into this category.
            \item Goal-based vs task-based.
        \end{itemize}
        
        \column{0.4\textwidth}
        \includegraphics<1>[width=\textwidth]{figures/agentic_demo.png}%
        \includegraphics<2>[width=\textwidth]{figures/agentic_demo2.png}%
    \end{columns}
\end{frame}

\begin{frame}
    \frametitle{Layer 3: Agentic Systems Cost}
    \framesubtitle{More expensive but potentially more valuable}
    
    \begin{itemize}
        \item Significantly more expensive than chat-based AI.
        \item Can consume hundreds of API calls per task.
        \item For agentic systems, subscriptions \emph{are} cheaper than pay-per-token.
        \item I was spending \$40/day on Claude Code with API tokens.
        \item Subscription gives you about 2hr coding session per day for \$20/month/user.
        \item Gemini CLI is ``free'' but risks going over free tier limit, and is less stable.
        \item In July 2025, Claude Code best in class.
        \item There are many further unknown costs of shifting your research toward ``hands-free coding''.
    \end{itemize}
\end{frame}

\begin{frame}
    \frametitle{Conclusions and Getting Started}
    
    \begin{columns}
        \column{0.5\textwidth}
        \begin{itemize}
            \item Use developer-level tools, not consumer interfaces.
            \item API access typically cheaper for research use.
            \item Capturing and owning content is critical -- use Otter for transcription.
            \item Export conversations and notes in formats that can be fed to LLMs.
            \item Get into the habit of recording information in exportable forms.
            \item Agentic systems are transformative but carry additional risks requiring careful consideration.
        \end{itemize}
        
        \column{0.5\textwidth}
        \begin{block}{To-do}
            \begin{itemize}
                \item Set up GitHub Copilot Pro (takes a couple of days to approve).
                \item Try AI Studio rather than ChatGPT.
                \item Get your PI to set up some tokens for your group to use (you only need ~O(\$10) to try it out).
                \item Try Google Gemini's Deep Research \textbf{button}.
                \item Test Claude Code/Gemini CLI for a month.
            \end{itemize}
        \end{block}
    \end{columns}
\end{frame}

\begin{frame}
    \frametitle{Miles Additions}
    
    \begin{columns}
        \column{0.5\textwidth}
        \begin{itemize}
            \item For comparing models across different tasks and capabilities:
            \item \tthref{lmarena.ai}
            \item Provides community-driven rankings based on real usage
            \item Updated regularly with new models and benchmarks
            \item Useful for selecting the right model for specific research needs
        \end{itemize}
        
        \column{0.5\textwidth}
        \includegraphics[width=\textwidth]{figures/lmarena_screenshot.png}
    \end{columns}
\end{frame}

\begin{frame}
    \frametitle{Claude Code \& OpenAI Codex}
    \framesubtitle{When you want to build something new}
    
    \begin{columns}
        \column{0.5\textwidth}
        \begin{itemize}
            \item If you want to make a new plot
            \item Add a new feature to a library
            \item Want to do a new analysis
            \item Claude Code: Interactive AI coding assistant
            \item OpenAI Codex: Code generation and completion
            \item Both excel at understanding context and generating working code
            \item Great for rapid prototyping and implementation
        \end{itemize}
        
        \column{0.5\textwidth}
        \includegraphics[width=\textwidth]{figures/claude_code_screenshot.png}
    \end{columns}
\end{frame}

\begin{frame}
    \frametitle{Cursor}
    \framesubtitle{AI-first code editor}
    
    \begin{columns}
        \column{0.5\textwidth}
        \begin{itemize}
            \item You should be using agentic systems as a first pass for a task
            \item VSCode-based editor with built-in AI
            \item Agent mode for autonomous coding
            \item Context-aware code generation
            \item Inline editing and completion
            \item Great for existing codebases
            \item Understands your project structure
        \end{itemize}
        
        \column{0.5\textwidth}
        \includegraphics[width=\textwidth]{figures/cursor_screenshot.png}
    \end{columns}
\end{frame}

\begin{frame}
    \frametitle{What percentage of your worktime do you spend?}
    
    \begin{itemize}
        \item Using AI tools for research tasks?
        \item How has this changed over the past year?
        \item What tasks do you find AI most/least helpful for?
        \item 0\% is too much, 100\% is too much
        \item Somewhere in between is good
        \item Most people probably don't spend enough time sharpening the saw
        \item We at the front of this lecture theatre probably spend too much time
    \end{itemize}
\end{frame}

\begin{frame}
    \frametitle{Getting Good at AI Tools}
    \framesubtitle{Practice and awareness}
    
    \begin{itemize}
        \item To get good at using these tools, the only way to do that is to use them a lot
        \item They're really good at lying
        \item They're really good at hallucinating and making things up
        \item Building up an awareness of when they are broken is very useful
        \item Practice helps you recognize patterns of failure
        \item Develop intuition for when outputs seem suspicious
        \item Always verify critical information independently
    \end{itemize}
\end{frame}

\begin{frame}
    \frametitle{Boris' Section: Agentic Tools}
    
    \begin{columns}
        \column{0.5\textwidth}
        \begin{itemize}
            \item Advanced AI systems that can act autonomously
            \item Take initiative and make decisions
            \item Multi-step problem solving
            \item Integration with external tools and APIs
            \item Examples: research assistants, code reviewers, data analysts
            \item Moving beyond simple question-answer interactions
            \item HAL/USACO is the agentic leaderboard for comparing performance
            \item \tthref{holistic-agentic-leaderboard.github.io}
            \item Doing things in-house with his group
            \item SWE-bench is a benchmark for how well an agent can solve GitHub issues automatically
        \end{itemize}
        
        \column{0.5\textwidth}
        \includegraphics[width=\textwidth]{figures/boris_agentic_screenshot.png}
    \end{columns}
\end{frame}

\begin{frame}
    \frametitle{Interfacing Agents with Your Own Documents or Code}
    
    \begin{itemize}
        \item Custom document integration for research workflows
        \item Code repository analysis and understanding
        \item Domain-specific knowledge incorporation
        \item RAG (Retrieval-Augmented Generation) systems
        \item Private knowledge bases and databases
        \item Maintaining context across large codebases
        \item Specialized tools for academic and research environments
    \end{itemize}
\end{frame}

\begin{frame}
    \frametitle{RAG means Retrieval Augmented Generation}
    
    \begin{columns}
        \column{0.5\textwidth}
        \begin{itemize}
            \item Combines retrieval of relevant documents with language generation
            \item First retrieve relevant information from your knowledge base
            \item Then generate responses using that retrieved context
            \item Reduces hallucinations by grounding in actual documents  
            \item Allows AI to work with your specific data and knowledge
            \item Examples: NotebookLM, ChatGPT with document uploads, custom RAG systems
            \item Essential for domain-specific or proprietary information
            \item<3> Created an assistant in the latest screenshot
        \end{itemize}
        
        \column{0.5\textwidth}
        \includegraphics<1>[width=\textwidth]{figures/vector_store_screenshot.png}%
        \includegraphics<2>[width=\textwidth]{figures/vector_store_screenshot2.png}%
        \includegraphics<3>[width=\textwidth]{figures/vector_store_screenshot3.png}%
    \end{columns}
\end{frame}

\begin{frame}
    \frametitle{Context Window of Current Models is Very Large}
    
    \begin{itemize}
        \item 2M tokens for Gemini 2.5 Pro
        \item Can process entire codebases or document collections
        \item Equivalent to ~1.5 million words or thousands of pages
        \item Enables comprehensive analysis without chunking
        \item Game-changer for research and development workflows
        \item No longer need complex RAG systems for many tasks
        \item Can maintain context across massive datasets
    \end{itemize}
\end{frame}

\begin{frame}
    \frametitle{Develop a Library in Python}
    
    \begin{itemize}
        \item Give the path of your codes/notebooks
        \item With a single click, turn all your files into a context file
        \item Which can be sent to an LLM
        \item Automatically processes multiple file types
        \item Handles code structure and documentation
        \item Optimizes for token efficiency
        \item Seamless integration with large context models
        \item Perfect for codebase analysis and documentation
        \item This is called contextmaker on PyPI
        \item Worked with the author of CAMB on CMB agent for this
    \end{itemize}
\end{frame}

\begin{frame}
    \frametitle{DENARIO Project}
    \framesubtitle{End-to-end paper writing system}
    
    \begin{columns}
        \column{0.5\textwidth}
        \begin{itemize}
            \item Automated academic paper writing system
            \item End-to-end research and writing pipeline
            \item Written dozens of papers from philosophy to astrophysics in the past months
            \item AI-driven research, analysis, and composition
            \item Cross-disciplinary capabilities
            \item Demonstrates the potential of agentic systems in academic research
        \end{itemize}
        
        \column{0.5\textwidth}
        \includegraphics<1>[width=\textwidth]{figures/denario_screenshot1.png}%
        \includegraphics<2>[width=\textwidth]{figures/denario_screenshot2.png}%
    \end{columns}
\end{frame}

\begin{frame}
    \frametitle{Additional Feedback: Outlier}
    
    \begin{itemize}
        \item Physics training company that gives free access to new LLMs
        \item Make you do a quick GCSE physics test for qualification
        \item Then provides free access to all the models
        \item Part of Scale AI's platform for AI model training
        \item Recruits experts to improve AI responses in physics
        \item \tthref{outlier.ai}
        \item Great way to get free access to cutting-edge models
    \end{itemize}
\end{frame}

\begin{frame}
    \frametitle{AI in Teaching: Durham University Research}
    
    \begin{columns}
        \column{0.48\textwidth}
        \begin{block}{Durham Study (GPT-3.5/4)}
            \begin{itemize}
                \item 42 exam papers, 593 questions (2018-2022)
                \item GPT-4: 49.4\% average
                \item GPT-3.5: 38.6\% average
                \item Used API to extract questions
                \item Human markers assessed responses
                \item Not blind marked
                \item Struggled with complex reasoning
            \end{itemize}
        \end{block}
        
        \column{0.48\textwidth}
        \begin{block}{My Experiment (o3)}
            \begin{itemize}
                \item Part II Cambridge astronomy exams
                \item Prompted as ``first class student''
                \item Full exam context and syllabus
                \item Transcribed onto exam scripts
                \item Marked blind alongside real scripts
                \item Exam not in training corpus
                \item \textbf{Result}: Best student in IoA history
            \end{itemize}
        \end{block}
    \end{columns}
    
    \vspace{0.3cm}
    \textbf{Key Difference}: Model generation (o3 vs GPT-4), prompting strategy, and exam complexity. Durham shows current limitations; Cambridge experiment shows frontier capabilities.
\end{frame}

\appendix

\end{document}
